% ch2.tex

\chapter{চাষাবাদ আর বর্ষপঞ্জি}\label{ch:2}

\epigraph{``\textit{
আমায় প্রশ্ন করে নীল ধ্রুব তারা,\\
আর কত কাল আমি রব দিশেহারা
}''}{\textit{সলিল চৌধুরী}}

% From the multitude of tribes living in a condition of barbarism, large states arose between 4000 and 1000 BC in the fertile plains of China, India, Mesopotamia and Egypt.
খ্রিস্টপূর্ব ৪০০০ থেকে ১০০০ অব্দের মধ্যে চীন, ভারত, মেসোপটেমিয়া ও মিশরের উর্বর সমভূমিতে বসবাসকারী বিভিন্ন জনগোষ্ঠী ধীরে ধীরে গড়ে তোলে বড় বড় রাষ্ট্র ও সভ্যতা।
% They represented a higher stage of culture, of which written documents have been transmitted to posterity.
তারা সংস্কৃতি ও সভ্যতার দিক থেকে আরও উন্নত পর্যায়ে পৌঁছেছিল। তাদের সেই অগ্রগতির প্রমাণ আজও পাওয়া যায় পরবর্তী প্রজন্মের জন্য রেখে যাওয়া বিভিন্ন লিখিত দলিল ও নথিপত্রে।
% The formerly independent peasant communities and town-ships, with their own chieftains and urban kings, their local gods and forms of worship, merged into larger political units.
আগে যেসব কৃষকসমাজ ও নগরকেন্দ্রিক জনপদ নিজেদের মতো স্বাধীনভাবে পরিচালিত হতো, যাদের ছিল নিজস্ব সর্দার বা নগরশাসক, স্থানীয় দেবদেবী ও উপাসনার রীতি, সেগুলো ধীরে ধীরে বৃহত্তর রাজনৈতিক ব্যবস্থার সঙ্গে একীভূত হয়ে যায়।
% The extraordinary fertility of the alluvial silt,  which yielded a surplus product, provided a living for a separate class of rulers and officiate.
উর্বর পলিমাটিতে প্রচুর ফসল ফলত। প্রয়োজন মেটানোর পর যে বাড়তি ফসল থেকে যেত, তা শাসক ও সরকারি কর্মকর্তাদের মতো পৃথক একটি শ্রেণির ভরণপোষণের ব্যবস্থা করত।
% This organization first arose from the need for the centralized regulation of the water.
পানির প্রবাহ ও বণ্টন কেন্দ্রীয়ভাবে নিয়ন্ত্রণের প্রয়োজন থেকেই প্রথম এই প্রশাসনিক ব্যবস্থার উদ্ভব ঘটে।
% The large rivers irrigating these plains the Nile, the Euphrates, the Hwang Ho filled their beds with silt, overflowed in certain months and flooded the fields, devastating or fertilizing them, or at times excavating a new bed.
এই সমভূমিগুলোর ওপর দিয়ে বয়ে যাওয়া নীল নদ, ইউফ্রেতিস ও হোয়াং হোর মতো বড় বড় নদী পলিমাটি জমিয়ে নিজেদের তলদেশ ভরাট করে ফেলত। বছরের নির্দিষ্ট সময়ে নদীগুলো দু'কূল ছাপিয়ে প্লাবিত করত আশপাশের কৃষিজমি। কখনো সেই বন্যায় ফসল ও জমির ক্ষতি হতো, আবার কখনো পলিমাটি ছড়িয়ে জমি হয়ে উঠত আরও উর্বর। কোনো কোনো সময় নদী তার পুরোনো গতিপথ ছেড়ে নতুন পথও তৈরি করত।
% The water had to be directed continuously and checked by dykes, by a deepening of the rivers, or by the digging of canals.
পানির প্রবাহ ঠিক রাখতে এবং তা নিয়ন্ত্রণে রাখতে নিয়মিত বাঁধ নির্মাণ, নদীর তলদেশ গভীর করা কিংবা নতুন খাল খননের প্রয়োজন হতো।
% Such control could not be left to the individual districts with their often conflicting interests. Centralized regulation was necessary, and only a strong central authority could guarantee that local interests would not prevail over general interests. Only then were fertility and prosperity assured.
বিভিন্ন অঞ্চলের স্বার্থ প্রায়ই পরস্পরের সঙ্গে সাংঘর্ষিক হতো। তাই পানির এই ব্যবস্থাপনার দায়িত্ব আলাদা আলাদা অঞ্চলের হাতে ছেড়ে দেওয়া সম্ভব ছিল না। প্রয়োজন ছিল কেন্দ্রীয় নিয়ন্ত্রণের, আর একমাত্র শক্তিশালী কেন্দ্রীয় শাসনব্যবস্থাই নিশ্চিত করতে পারত যে সামগ্রিক স্বার্থের চেয়ে কোনো অঞ্চলের নিজস্ব স্বার্থ বড় হয়ে উঠবে না। এভাবেই নিশ্চিত করা সম্ভব হতো জমির উর্বরতা ও সমৃদ্ধি।
% But when the country fell asunder into petty warring principalities and the dykes and canals were neglected, the soil dried up or was flooded, and the people starved: the wrath of the gods then lay on the land.
কিন্তু দেশ যখন ভেঙে ছোট ছোট বিবদমান রাজ্যে বিভক্ত হয়ে পড়ত এবং সেগুলোর মধ্যে যুদ্ধ লেগেই থাকত, তখন বাঁধ ও খালগুলোর যথাযথ রক্ষণাবেক্ষণ হতো না। ফলে কখনো জমি শুকিয়ে ফসল নষ্ট হতো, আবার কখনো বন্যায় ডুবে যেত। শেষ পর্যন্ত মানুষ দুর্ভিক্ষের মুখে পড়ত। তখন এই বিপর্যয়কে মনে করা হতো দেবতাদের অভিশাপ বা রোষের ফল।



% A strong central power was needed, secondly, for the defence of the fertile plains against the war-like inhabitants of the adjacent mountains or deserts.
শক্তিশালী কেন্দ্রীয় শাসনব্যবস্থার প্রয়োজন ছিল আরও একটি কারণে—পাশের পাহাড়ি ও মরু অঞ্চলে বসবাসকারী যুদ্ধপ্রবণ জনগোষ্ঠীর আক্রমণ থেকে উর্বর সমভূমিগুলোকে রক্ষা করার জন্য।
% These people could find only a meagre living on their own land, so they made it their business to plunder and to force tribute from their prosperous neighbours.
নিজেদের ভূমিতে জীবিকা নির্বাহ করাই ছিল তাদের জন্য কঠিন। তাই তারা আশপাশের সমৃদ্ধ জনপদে লুটতরাজ চালাত এবং সেখানকার মানুষের কাছ থেকে জোর করে কর বা খাজনা আদায় করত।
% A division of tasks was necessary; a caste of warriors developed who, with their chief as king, became a ruling class controlling the surplus product of the farmers.
এভাবে কাজের দায়িত্ব ভাগ করে দেওয়ার প্রয়োজন থেকে একটি পৃথক যোদ্ধাশ্রেণির উদ্ভব ঘটে। তাদের নেতা রাজা হিসেবে প্রতিষ্ঠিত হন এবং এই যোদ্ধাশ্রেণিই ধীরে ধীরে শাসকগোষ্ঠীতে পরিণত হয়ে কৃষকদের উৎপাদিত বাড়তি ফসলের ওপর নিয়ন্ত্রণ প্রতিষ্ঠা করে।
% Or the roving neighbours turned from marauders into conquerors and settled as a ruling aristocracy in the midst of the farmers, protecting them from other aggressors. In either case the result was a strong centralized state power.
কখনো আশপাশের যাযাবর জনগোষ্ঠী লুটপাটের পরিবর্তে ভূখণ্ড দখল করে সেখানে স্থায়ীভাবে বসতি গড়ে তুলত। এভাবে তারা কৃষকদের মাঝে শাসক অভিজাত শ্রেণি হিসেবে প্রতিষ্ঠিত হতো এবং অন্য আক্রমণকারীদের হাত থেকে কৃষকদের রক্ষা করত। উভয় ক্ষেত্রেই শেষ পর্যন্ত গড়ে উঠত একটি শক্তিশালী কেন্দ্রীয় শাসনব্যবস্থা।



% This story, repeated from time to time, represents in brief outline the history of these lands. Over and over again barbarian peoples invaded and subjected them, sometimes remaining as a thin upper layer, like the Manchus in China, the Mongols in India, the Hyksos in Egypt, and sometimes as an entire people mixing with or replacing the former inhabitants, like the Aryans in India and the Semites in Mesopotamia.
এই ভূখণ্ডগুলোর ইতিহাস পর্যালোচনা করলে বারবার একই ঐতিহাসিক ধারার পুনরাবৃত্তি চোখে পড়ে। বিভিন্ন সময়ে নানা যুদ্ধবাজ জাতিগোষ্ঠী এসব ভূখণ্ডে আক্রমণ চালিয়েছে। কখনো তারা বিজিত অঞ্চলের ওপর শাসকগোষ্ঠীর একটি ক্ষুদ্র উপরিস্তর হিসেবে নিজেদের প্রতিষ্ঠিত করেছে, অথচ মূল শাসনক্ষমতা বজায় রেখেছে নিজস্ব মাতৃভূমিতে—যেমন চীনে মাঞ্চুরা, ভারতে মোগলরা এবং মিশরে হিকসসরা। আবার কখনো তারা স্থানীয় জনগোষ্ঠীর সঙ্গে মিশে গিয়ে তাদেরই অংশ হয়ে উঠেছে, যেমন ভারতে আর্যরা এবং মেসোপটেমিয়ায় সেমিটিক জনগোষ্ঠী। কোনো কোনো ক্ষেত্রে স্থানীয় জনগোষ্ঠীকে সম্পূর্ণরূপে বিতাড়িত বা উৎখাতও করেছে।
% Though much culture was destroyed in the conquest, the invaders afterwards adopted and assimilated the existing higher state of civilization, and often imparted fresh vigour to it. After some generations, having lost their savage strength, the conquerors themselves became an object of attack by new aggressors.
আগ্রাসনের ফলে অনেক সময় স্থানীয় সংস্কৃতি ও ঐতিহ্য বিলুপ্ত হলেও, বিজয়ী জনগোষ্ঠী পরবর্তীকালে বিজিত অঞ্চলের অপেক্ষাকৃত উন্নত সভ্যতা ও সংস্কৃতিকে গ্রহণ করে নিজেদের মধ্যে আত্মস্থ করেছে। এর মধ্য দিয়ে তাদের সমাজে নতুন প্রাণশক্তি ও উদ্যমের সঞ্চার হয়েছে। তবে কয়েক প্রজন্মের ব্যবধানে বিজয়ী শাসকগোষ্ঠীও নিজেদের আগের সেই দুর্ধর্ষ শক্তি ও সামরিক দাপট হারিয়ে ফেলত। একসময় তারাই নতুন আক্রমণকারীদের আগ্রাসনের শিকারে পরিণত হতো।




In such empires the prince was the legislator., the chief justice and the head of a government of officials, who as leaders of the civil departments constituted a second ruling class alongside the military. Usually it consisted ofthe priests, who had been the local intellectual leaders, and now were organized into an official hierarchy. The priesthood held the spiritual leadership of state and society in its hands. It held the theoretical and general knowledge needed in the process of production; this was the source of its prestige and its social power. Where agriculture was the main occupation, knowledge of the calendar and of the seasons was their domain. Religion too in those times one with state and society was centralized; the local deities of the chief townships were conjoined into a pantheon under an upper god, whereas the local gods of smaller or conquered sites were degraded to a lower order of spirits or were incorporated into other gods. Thus in ancient Babylonia the goddesses Innina, Nisaba and Nana were united with the later Ishtar.

The fact that Borsippa became a suburb of the rapidly growing city of Babylon found its theological expression in its city god Nabo becoming the son ofMarduk. Whereas the deities of the oldest centres of culture Eridu, Uruk and Nippur always remained the most highly venerated gods Ea, Ami and Enlil (mostly called Bel, i.e. 'Lord'), after 2000 BG the political hegemony of Babylon made its local god Marduk the senior god of the pantheon. When in later centuries Assyria acquired the overlordship of the Mesopotamian world, Ashur took this place. When combined into a hierarchy, the gods lost some of their local character, and their character as personifications of natural powers became more prominent. The deities Sin of Ur and Shamash of Sippara in later times were always venerated as the moon god and the sun god.

এমন সাম্রাজ্যে রাজা ছিলেন আইনপ্রণেতা, প্রধান বিচারক এবং কর্মকর্তাদের নিয়ে গঠিত সরকারের প্রধান। বেসামরিক বিভাগগুলোর নেতা হিসেবে এই কর্মকর্তারা সামরিক শ্রেণির পাশাপাশি দ্বিতীয় একটি শাসকশ্রেণি গঠন করতেন। সাধারণত এই শ্রেণি গঠিত হতো পুরোহিতদের নিয়ে। তাঁরা আগে স্থানীয় বুদ্ধিবৃত্তিক নেতা ছিলেন, এখন একটি দাপ্তরিক স্তরবিন্যাসে সংগঠিত হলেন। রাষ্ট্র ও সমাজের আধ্যাত্মিক নেতৃত্ব ছিল পুরোহিততন্ত্রের হাতে। উৎপাদনপ্রক্রিয়ার জন্য প্রয়োজনীয় তাত্ত্বিক ও সাধারণ জ্ঞান তাঁদের কাছেই ছিল। এটাই ছিল তাঁদের মর্যাদা ও সামাজিক ক্ষমতার উৎস। যেখানে কৃষিই ছিল প্রধান পেশা, সেখানে পঞ্জিকা ও ঋতুর জ্ঞান ছিল তাঁদের এখতিয়ারে।

সে যুগে ধর্মও রাষ্ট্র ও সমাজের সঙ্গে একাকার ছিল, এবং ধর্মও কেন্দ্রীভূত হলো। প্রধান নগরগুলোর স্থানীয় দেবতারা একজন প্রধান দেবতার অধীনে একটি দেবমণ্ডলীতে একত্রিত হলেন। অন্যদিকে ছোট বা বিজিত স্থানগুলোর স্থানীয় দেবতারা নিম্নস্তরের আত্মায় নেমে গেলেন, কিংবা অন্য দেবতাদের মধ্যে মিশে গেলেন। যেমন প্রাচীন ব্যাবিলোনিয়ায় দেবী ইন্নিনা, নিসাবা ও নানা পরবর্তীকালের ইশতারের সঙ্গে এক হয়ে যান। দ্রুত বেড়ে ওঠা ব্যাবিলন নগরীর শহরতলিতে পরিণত হয় বোরসিপ্পা। এর ধর্মতাত্ত্বিক প্রকাশ ঘটে এভাবে যে বোরসিপ্পার নগরদেবতা নাবো মারদুকের পুত্র হয়ে যান। সংস্কৃতির প্রাচীনতম কেন্দ্র এরিদু, উরুক ও নিপ্পুরের দেবতারা, অর্থাৎ এয়া, আনু ও এনলিল (যাঁকে প্রায়ই বেল, অর্থাৎ "প্রভু" বলা হতো), সবসময়ই সর্বাধিক পূজিত দেবতা ছিলেন। তবু খ্রিস্টপূর্ব ২০০০ সালের পর ব্যাবিলনের রাজনৈতিক আধিপত্যের ফলে তার স্থানীয় দেবতা মারদুক দেবমণ্ডলীর প্রধান দেবতা হয়ে ওঠেন। পরের শতাব্দীগুলোতে আসিরিয়া যখন মেসোপটেমীয় জগতের ওপর কর্তৃত্ব লাভ করে, তখন আশুর এই স্থান নেন। স্তরবিন্যাসে একত্রিত হওয়ার পর দেবতারা তাঁদের স্থানীয় চরিত্র কিছুটা হারান, আর প্রাকৃতিক শক্তির মূর্ত রূপ হিসেবে তাঁদের চরিত্র বেশি প্রকট হয়ে ওঠে। যেমন উরের দেবতা সিন এবং সিপ্পারার দেবতা শামাশ পরবর্তীকালে সবসময় চন্দ্রদেব ও সূর্যদেব হিসেবে পূজিত হয়েছেন।

এমন একটি গোষ্ঠীর উত্থান ঘটল, যাদের শাসকশ্রেণি হিসেবে আর কঠোর পরিশ্রম করে জীবন নিরাপদ রাখতে হতো না। এর ফলে অস্তিত্বের নতুন পরিস্থিতি তৈরি হলো। সমাজকাঠামো আরও জটিল হলো, আর আধ্যাত্মিক নেতৃত্বের জন্য উচ্চতর যোগ্যতার প্রয়োজন হলো এবং তার দাবিও বিস্তৃত হলো। ব্যবসা-বাণিজ্য নতুন বস্তুগত ও আত্মিক চাহিদার জন্ম দিল। রাজা ও অভিজাতদের সম্পদ ও বিলাসিতার সঙ্গে এল শিল্প ও বিজ্ঞানের প্রতি অনুরাগ। এভাবে পৃথিবীতে প্রথমবারের মতো নতুন সমাজকাঠামোর পাশাপাশি সংস্কৃতির উচ্চতর রূপের উদ্ভব হলো, যা প্রাগৈতিহাসিক বর্বরতার সবচেয়ে উন্নত রূপকেও ছাড়িয়ে গেল। শুরু হলো সভ্যতার যুগ।

এর সঙ্গে জ্যোতির্বিজ্ঞানেরও উচ্চতর বিকাশ ঘটল। এর উৎস ছিল সরাসরি কালগণনার প্রয়োজন, বিশেষ করে চান্দ্র পঞ্জিকাকে সৌর বছরের সঙ্গে খাপ খাওয়ানোর সমস্যা।

চাঁদের একটি পর্যায়কাল গড়ে ২৯.৫৩০৫৯ দিন। একটি সৌর বছর ৩৬৫.২৪২২০ দিন, যা ১২টি চান্দ্র পর্যায়কালের (মোট ৩৫৪.৩৬৭১ দিন) চেয়ে ১১ দিন বেশি। তিন বছর পর চান্দ্র পঞ্জিকা সূর্যের অগ্রগতি থেকে ৩৩ দিন পিছিয়ে পড়ে। সূর্যের সঙ্গে সামঞ্জস্য রাখতে প্রতি তৃতীয় বছরে, কখনো আরও ঘন ঘন, একটি বাড়তি মাস যোগ করতে হতো, যাতে বছরে ১২টির বদলে ১৩টি মাস হয়। পঞ্জিকার সমস্যাটি হলো এমন একটি বড় পর্যায়কাল খুঁজে বের করা, যা মাস ও বছর উভয়েরই সাধারণ গুণিতক। তাহলে এই পর্যায়কাল শেষে সূর্য ও চাঁদ আবার একই পারস্পরিক অবস্থানে ফিরে আসবে। অবশ্য নিখুঁত কোনো সাধারণ গুণিতক নেই, তবে কমবেশি সন্তোষজনক আসন্ন মান পাওয়া যায়।

সৌর ও চান্দ্র পর্যায়কাল সম্পর্কে আমাদের নির্ভুল আধুনিক জ্ঞান দিয়ে এমন আসন্ন পঞ্জিকা-পর্যায়কাল তাত্ত্বিকভাবে বের করা যায়। এর জন্য দুটির অনুপাতকে একটি অবিচ্ছিন্ন ভগ্নাংশে (continued fraction) রূপান্তর করে তার ক্রমিক আসন্ন মানগুলো লিখতে হয়। এভাবে আমরা ৩৭/৩ ও ৯৯/৮ পাই। প্রথমটি দেখায় যে ৮ বছর প্রায় ৯৯ মাসের সমান (যথাক্রমে ২৯২১.৯৪ ও ২৯২৩.৫৩ দিন)। তাই এই ৮ বছরের মধ্যে ৩টিতে ১৩তম মাস থাকতে হবে, আর ৫টিতে থাকবে ১২টি মাস। তবু এই আসন্ন মান খুব ভালো নয়। মাত্র ২৪ বছর পরই চান্দ্র তারিখ সৌর ঋতুর তুলনায় ৫ দিন পিছিয়ে পড়বে। এর চেয়ে অনেক বেশি নির্ভুল হলো ১৯ বছরের একটি পর্যায়কাল, যাতে ২৩৫টি মাস থাকে (যথাক্রমে ৬৯৩৯.৬০ ও ৬৯৩৯.৬৯ দিন)। এখানে ১৩তম মাস সাতবার যোগ করতে হয়। অবশ্যই আদিতম যুগে মানুষের পর্যায়কাল সম্পর্কে এই নিখুঁত জ্ঞান ছিল না। তাদের জন্য একটি ভালো পঞ্জিকা-পর্যায়কাল খুঁজে পাওয়া ছিল কঠিন এক ব্যবহারিক সমস্যা। সৌর ও চান্দ্র গণনাকে খাপ খাওয়ানোর শ্রমসাধ্য প্রক্রিয়ার মাধ্যমেই কেবল এর সমাধান সম্ভব ছিল। তাই এই সমস্যা আকাশের ঘটনাবলিকে আরও যত্নের সঙ্গে পর্যবেক্ষণ করার চালিকাশক্তি হিসেবে কাজ করেছিল।

প্রশ্ন উঠতে পারে, এত নির্ভুলতার প্রয়োজন কেন ছিল? এটা তো কৃষির প্রয়োজনের চেয়ে অনেক বেশি, আর আবহাওয়ার ওঠানামার কারণে কৃষিকাজ এমনিতেই বেশ অনিয়মিত। তবে আমাদের মনে রাখতে হবে, সে যুগে কৃষিকাজের সঙ্গে ধর্মীয় আচার ও উৎসব জড়িত থাকত। অন্য সব বড় ও গুরুত্বপূর্ণ সামাজিক ঘটনার মতো কৃষি-উৎসবগুলোও একই সঙ্গে ছিল ধর্মীয় অনুষ্ঠান। প্রধান প্রাকৃতিক ও সামাজিক শক্তিগুলোর প্রতিনিধি হিসেবে দেবতারা মানুষের জীবনে অংশ নিতেন। সামাজিকভাবে যা প্রয়োজনীয় বা উপযুক্ত ছিল, তা দেবতাদের আদেশে পরিণত হতো এবং আচার-অনুষ্ঠানে কঠোরভাবে নির্দিষ্ট করে দেওয়া হতো। প্রকৃতিগতভাবে যা একটি নির্দিষ্ট ঋতুতে ঘটত, যেমন ফসল ঘরে তোলার উৎসব, ধর্মীয় অনুষ্ঠান হিসেবে তা একটি নির্দিষ্ট তারিখে বেঁধে দেওয়া হতো, যেমন চাঁদের কোনো একটি কলায়। দেবতাদের সেবায় কোনো অবহেলার সুযোগ ছিল না, সেখানে আচার-অনুষ্ঠান নিখুঁতভাবে পালন করতে হতো। পঞ্জিকা মূলত ছিল আচার-অনুষ্ঠানের কালানুক্রমিক বিন্যাস। এই কারণে পঞ্জিকা দায়িত্বপ্রাপ্ত কর্মকর্তাদের, বিশেষ করে পুরোহিতদের, নিরন্তর যত্নের বিষয় হয়ে উঠেছিল। আবার শুভ সময়ের জ্ঞানে তাঁদের একচেটিয়া অধিকার থাকায় পঞ্জিকা একই সঙ্গে তাঁদের সামাজিক ক্ষমতারও উৎস ছিল।

এ ধরনের ধর্মীয়-কৃষিভিত্তিক চর্চা থেকে কীভাবে চান্দ্র-সৌর পঞ্জিকা প্রতিষ্ঠিত হলো, তার কয়েকটি চমকপ্রদ দৃষ্টান্ত আমাদের কাছে এসে পৌঁছেছে। ইহুদি পঞ্জিকার উৎপত্তি মরুভূমিতে, তাই এটি চাঁদের ওপর ভিত্তি করে তৈরি হয়েছিল। কিন্তু ইহুদিরা কেনান দেশে পৌঁছানোর পর কৃষিই প্রধান পেশা হয়ে উঠল, আর পঞ্জিকাকে সূর্যের সঙ্গে খাপ খাওয়াতে হলো। বসন্তকালে মাসোথ ছিল ফসল তোলার উৎসব। এতে যবের প্রথম আঁটিগুলো নিবেদন করা হতো, আর প্রথম শস্য দিয়ে খামিরবিহীন রুটি তৈরি করা হতো। এই উৎসব যাযাবরদের পেসাখ উৎসবের সঙ্গে মিশে যায়। পেসাখে বসন্তের পূর্ণিমায় সদ্যোজাত মেষশাবক জিহোভার উদ্দেশে উৎসর্গ করা হতো। এভাবে উৎসবটি প্রথম মাস নিসানের পূর্ণিমায় স্থির হলো।

কীভাবে এটি করা হতো, তা গিনৎসেলের কালগণনা-বিষয়ক বিখ্যাত পাঠ্যবইয়ে পড়া যায়: "শেষ মাসের শেষ দিকে পুরোহিতেরা মাঠে গিয়ে ফসলের অবস্থা পরিদর্শন করতেন, পরের দুই সপ্তাহে যব পাকবে বলে আশা করা যায় কি না তা দেখতে। যদি দেখতেন যে পাকবে, তাহলে পরবর্তী অমাবস্যা থেকে শুরু হওয়া মাসেই মাসোথ-পেসাখ নির্ধারণ করা হতো। কিন্তু যদি পাকার আশা না থাকত, তাহলে উৎসবের মাস এক চান্দ্র পর্যায়কাল পিছিয়ে দেওয়া হতো। এর ভিত্তিতেই অন্যান্য উৎসবের দিন নির্ধারিত হতো।"

ইসরায়েলিরা যতদিন প্যালেস্টাইনে একসঙ্গে বাস করত, ততদিন এই অভিজ্ঞতাভিত্তিক পদ্ধতি তাদের কাজে লেগেছিল। পরে তারা ছড়িয়ে পড়লে এটি আর ব্যবহার করা গেল না। তবে ততদিনে জ্যোতির্বিজ্ঞান এতটা এগিয়ে গিয়েছিল যে তারা প্রতিবেশীদের কাছ থেকে পঞ্জিকা-পর্যায়কালের জ্ঞান গ্রহণ করতে পারল। তখন ১৯ বছরের পর্যায়কাল ইহুদি পঞ্জিকার ভিত্তি হয়ে উঠল। এতে ১২ মাসের ১২টি বছর আর ১৩ মাসের ৭টি বছর একটি নির্দিষ্ট ক্রমে পালাক্রমে আসে।

আরেকটি দৃষ্টান্ত হলো মুহাম্মদের পূর্ববর্তী আরবের ঘটনাপ্রবাহ। পবিত্র মাসগুলোতে রক্তের প্রতিশোধ স্থগিত রাখা হতো এবং কাফেলাগুলো নিরাপদে চলাচল করতে পারত, কারণ অর্থনৈতিক জীবনের জন্য এটা দরকার ছিল। গোত্রগুলোর মধ্যে অন্তহীন বিবাদ ছিল খুবই স্বাভাবিক ব্যাপার। গোত্রের সদস্যদের মধ্যে সংহতির একটি আদিম বিচারিক রূপ হিসেবে রক্তের প্রতিশোধও প্রয়োজনীয় ছিল। কিন্তু এর সীমাহীন দাপট চললে বেচাকেনা আর খাদ্যের জোগান অসম্ভব হয়ে পড়ত। মাসউদী বলেছেন: "সফর মাসের নাম এসেছে ইয়েমেনের বাজার থেকে… এখানে আরবরা তাদের শস্য কিনত, আর যে তা করতে ভুলত সে না খেয়ে মারা যেত।"

পবিত্র মাস ছিল বাজারের মাস। চারদিক থেকে কাফেলাগুলো বড় বড় বাজার-শহরে যেত, বিশেষ করে মক্কায়। মক্কা তার অতুলনীয় শীতল কূপ জমজমের জন্য বিখ্যাত ছিল, যা আরবে তীর্থকেন্দ্র হিসেবে মক্কার গুরুত্বের উৎস। এখানে গোত্রগুলো মিলিত হতো ও আলাপ-আলোচনা করত, আর এই মিলনস্থলই রাজনৈতিক ও ধর্মীয় কেন্দ্র এবং রাজধানীতে পরিণত হলো। তাই পবিত্র মাসটি সবসময় বছরের একই সময়ে পড়া জরুরি ছিল, যখন ফসল তোলা হয়ে গেছে এবং পণ্য হাতে আছে।

পরবর্তীকালের মুসলিম পণ্ডিত আল-বিরুনী লিখেছেন: "জাহেলি যুগে আরবরা তাদের মাসগুলোকে সেভাবেই গণনা করত যেভাবে মুসলমানরা এখন করে, আর তাদের হজ বছরের চারটি ঋতু জুড়ে ঘুরে বেড়াত। তবে পরে তারা স্থির করল যে হজ এমন সময়ে হবে যখন তাদের পণ্য, চামড়া ও ফল বাজারে বিক্রির জন্য প্রস্তুত থাকে। তাই তারা একে অচল করতে চাইল, যাতে এটি সবসময় সবচেয়ে প্রাচুর্যের ঋতুতে পড়ে। এজন্য হিজরতের ২০০ বছর আগে তারা ইহুদিদের কাছ থেকে অধিমাস যোগ করার পদ্ধতি শিখল। ইহুদিদের মতো করেই তারা এটি প্রয়োগ করল: তাদের বছর আর সৌর বছরের পার্থক্য যখন বেড়ে এক মাসের সমান হতো, তখন তা বছরের মাসগুলোর সঙ্গে যোগ করত। হজের অনুষ্ঠান শেষে কালাম্মাসরা (এই দায়িত্বপ্রাপ্ত একটি গোত্রের শেখরা) সামনে এসে জনতার উদ্দেশে বক্তৃতা দিতেন। তাঁরা পরের মাসটিকে চলতি মাসের নাম দিয়ে একটি মাস যোগ করতেন। জনতা হাততালি দিয়ে কালাম্মাসদের সিদ্ধান্তে সম্মতি জানাত। এই প্রক্রিয়াকে তারা বলত নাসি, অর্থাৎ স্থানান্তর, কারণ প্রতি দ্বিতীয় বা তৃতীয় বছরে বছরের শুরু সরে যেত… মানাজিলের উদয় ও অস্ত দেখে তারা এটি বিচার করতে পারত। মক্কা থেকে মদিনায় নবীর হিজরত পর্যন্ত এই ব্যবস্থাই চালু ছিল।"

হিজরতের পর নবম বছরে মুহাম্মদ এই স্থানান্তর নিষিদ্ধ করেন। হয়তো কালাম্মাসদের দায়িত্ব কেড়ে নিয়ে তাঁদের আধ্যাত্মিক ক্ষমতা ভেঙে দেওয়া এবং ইহুদিদের থেকে নিজেকে আরও স্পষ্টভাবে আলাদা করাই ছিল এর উদ্দেশ্য। তাই এখন মুসলিম পঞ্জিকা ৩৫৪ দিনের চান্দ্র বছরের ওপর ভিত্তি করে চলে। প্রতিটি বছরে ১২টি চান্দ্র মাস থাকে, আর ৩৩ বছরে এগুলো সব ঋতু ঘুরে আসে। এখানে আমরা এমন একটি পঞ্জিকা পাই যা আনুষ্ঠানিক বিধানের ফলে সামাজিক চর্চা থেকে বিচ্ছিন্ন হয়ে গেছে। এটি আদিম বেদুইন জীবনধারার এক প্রস্তরীভূত ধর্মীয় ঐতিহ্যে পরিণত হয়েছে।

এটি দেখাতে পারে যে ভিন্ন ভিন্ন অর্থনৈতিক পরিস্থিতি ও রাজনৈতিক ইতিহাসের ফলে বিভিন্ন জাতির মধ্যে পঞ্জিকা ও জ্যোতির্বিজ্ঞান কীভাবে ভিন্ন ভিন্ন পথে বিকশিত হয়েছে।
